% ============================================================
% adicoes_memoria_tcc.tex
%
% Todas as adições e substituições necessárias para integrar
% a análise de memória ao documento do TCC.
%
% Cada bloco indica:
%   [LOCALIZAÇÃO] — onde encontrar o trecho original
%   [AÇÃO]        — o que fazer (substituir / adicionar / remover)
%   [ORIGINAL]    — trecho exato a ser localizado no documento
%   [NOVO]        — código LaTeX que deve ficar no lugar
%
% Ordem sugerida de aplicação: de baixo para cima no documento
% evita que mudanças anteriores desorientem a busca.
% ============================================================


% ============================================================
% BLOCO 1 — Resumo (Abstract em português)
%
% LOCALIZAÇÃO: parágrafo que lista as métricas avaliadas no Resumo.
% AÇÃO: Substituir o trecho que menciona as métricas por versão
%       que inclui "consumo de memória".
%
% ORIGINAL (localize esta frase):
%   "tamanho de chaves e assinaturas, tempo de execução"
%   (ou qualquer lista de métricas que NÃO inclua memória)
%
% NOVO: acrescente ", e consumo de memória (pilha e heap)"
%       ao final da lista de métricas. Exemplo:
% ============================================================
%
% Exemplo de como a frase deve ficar:
%   "...avaliando tamanho de chaves e assinaturas, tempo de
%    execução, variação de flags de segurança e consumo de
%    memória (\textit{stack} e \textit{heap})."
%
% *** Ajuste conforme o texto exato do seu Resumo. ***


% ============================================================
% BLOCO 2 — Abstract (em inglês)
%
% LOCALIZAÇÃO: parágrafo que lista as métricas no Abstract.
% AÇÃO: Mesmo que o Bloco 1 — adicionar "memory consumption
%       (stack and heap)" à lista de metrics.
%
% Exemplo de como a frase deve ficar:
%   "...evaluating key and signature sizes, execution time,
%    security flag variation, and memory consumption
%    (stack and heap)."
%
% *** Ajuste conforme o texto exato do seu Abstract. ***


% ============================================================
% BLOCO 3 — Tabela de Trabalhos Correlatos (tab:correlatos)
%
% LOCALIZAÇÃO: linha "\textbf{Este trabalho}" na tabela correlatos.
% AÇÃO: Na coluna "Métricas Principais" dessa linha, acrescentar
%       "consumo de memória" à lista existente.
%
% ORIGINAL (célula atual, aproximadamente):
%   Tempo de execução, variação de flags e níveis NIST
%
% NOVO (substitua apenas o conteúdo dessa célula):
% ============================================================
%
%   Tempo de exec., tamanho de chaves/assinaturas,
%   variação de flags, níveis NIST e consumo de
%   memória (\textit{stack}/\textit{heap})
%
% *** Localize a linha deste trabalho e atualize a célula. ***


% ============================================================
% BLOCO 4 — Seção 3.3 (Métricas de Avaliação)
%
% LOCALIZAÇÃO: lista de métricas dentro de \subsection ou
%   parágrafo equivalente na Seção 3. Deve haver itens como
%   "Tempo de Execução", "Tamanho de Chaves", etc.
% AÇÃO: Adicionar o item abaixo como 5º (ou último) item da lista.
%
% Se a seção usa \begin{itemize}...\end{itemize}, insira antes
% do \end{itemize}. Se usa parágrafos numerados, adicione ao final.
% ============================================================

% --- INSERIR antes do \end{itemize} ou ao final da lista de métricas ---
\item \textbf{Consumo de Memória.}
O uso de memória é avaliado por meio de duas técnicas complementares.
O consumo de \textit{stack} é medido com a técnica de
\textit{stack painting}: um \textit{buffer} sentinela de 4\,MB é
preenchido com o byte~\texttt{0xAB} antes da operação; ao término,
varredura do endereço mais baixo identifica quantos bytes foram
efetivamente escritos pela \textit{thread}.
O pico de \textit{heap} é obtido com o
\texttt{valgrind --tool=massif}, que registra instantâneos do campo
\texttt{mem\_heap\_B} durante toda a execução e retém o valor máximo.
Ambas as métricas são coletadas com a flag de otimização
\texttt{-O0} para preservar a estrutura real de chamadas.


% ============================================================
% BLOCO 5 — Seção 4 (Implementação)
%
% LOCALIZAÇÃO: final das subseções de implementação (após a
%   descrição dos benchmarks de tempo/tamanho), ainda dentro
%   da Seção 4.
% AÇÃO: Adicionar nova \subsection descrevendo a infraestrutura
%       de medição de memória.
% ============================================================

% --- INSERIR como nova subseção ao final da Seção 4 ---
\subsection{Medição do Consumo de Memória}
\label{sec:impl-memoria}

A medição de memória é realizada por dois programas distintos,
compilados com \texttt{-O0 -g} para preservar o comportamento
de alocação sem otimizações do compilador.

\textbf{Stack painting (\texttt{memory\_benchmark.c}).}
Para cada operação (\textit{keygen}, \textit{encap}/\textit{sign},
\textit{decap}/\textit{verify}), uma \textit{thread} POSIX é criada
com \texttt{pthread\_attr\_setstack()} apontando para um
\textit{buffer} de 4\,MB previamente pintado com o byte
sentinela~\texttt{0xAB}.
Após o retorno da operação, o \textit{buffer} é varrido do endereço
mais baixo em direção ao mais alto; o primeiro byte diferente
de~\texttt{0xAB} delimita a região de \textit{stack} não utilizada,
e o uso real é calculado como o complemento.
O resultado reportado é a mediana de cinco execuções após dois
aquecimentos (\textit{warm-up}).

\textbf{Valgrind massif (\texttt{massif\_single.c}).}
Um programa auxiliar executa exatamente \emph{uma} chamada da
operação desejada por invocação, permitindo que o
\texttt{valgrind --tool=massif} perfile o \textit{heap} daquela
operação isolada.
O \textit{script} \texttt{run\_massif\_all.sh} automatiza a execução
para todos os algoritmos e operações; o \textit{script}
\texttt{parse\_massif.py} extrai o valor de pico de
\texttt{mem\_heap\_B} de cada arquivo \texttt{.massif} gerado e
consolida os resultados em CSV\@.

Como a versão do \texttt{liboqs} utilizada ainda adota os nomes
internos originais (p.\,ex., \texttt{Falcon-512} em vez de
\texttt{FN-DSA-512}), o programa \texttt{massif\_single.c} inclui
uma tabela de mapeamento que traduz os nomes FIPS/NIST para os
identificadores aceitos pela API\@.
As medições de \textit{stack} para as variantes
SLH-DSA-SHAKE não estão disponíveis nesta versão do
\textit{benchmark} e estão sinalizadas como indisponíveis nas
tabelas correspondentes.


% ============================================================
% BLOCO 6 — Seção 5, parágrafo introdutório
%
% LOCALIZAÇÃO: primeiro parágrafo da Seção 5 (Resultados),
%   que provavelmente menciona "quatro dimensões" ou lista
%   as subseções que virão.
% AÇÃO: Substituir "quatro dimensões" por "cinco dimensões"
%   e acrescentar a referência à subseção de memória.
%
% ORIGINAL (localize e substitua):
%   "quatro dimensões"  →  "cinco dimensões"
%
% Se houver uma lista das subseções, acrescente ao final:
%   "e consumo de memória (Seção~\ref{sec:res-memoria})"
%
% Exemplo de como o parágrafo pode ficar:
% ============================================================
%
% Esta seção apresenta os resultados experimentais organizados
% em cinco dimensões: tamanho de chaves e assinaturas
% (Seção~\ref{sec:res-tamanho}), tempo de execução
% (Seção~\ref{sec:res-tempo}), variação de \textit{flags} de
% segurança (Seção~\ref{sec:res-flags}), comparação entre
% níveis de segurança (Seção~\ref{sec:res-niveis}) e consumo
% de memória (Seção~\ref{sec:res-memoria}).
%
% *** Ajuste os \ref{} conforme os labels reais do documento. ***


% ============================================================
% BLOCO 7 — Nova Subseção 5.5: Consumo de Memória
%
% LOCALIZAÇÃO: após a última subseção atual da Seção 5
%   (provavelmente 5.4 — Comparação por Nível de Segurança).
% AÇÃO: Inserir todo o bloco abaixo.
% ============================================================

% --- INSERIR após a última subseção da Seção 5 ---
\subsection{Consumo de Memória}
\label{sec:res-memoria}

O consumo de memória foi avaliado por duas técnicas complementares:
\textit{stack painting} para medir o uso de pilha e
\texttt{valgrind --tool=massif} para medir o pico de \textit{heap}.
Ambas as medições foram realizadas com compilação \texttt{-O0},
sem otimizações, para refletir o comportamento real de alocação
de cada algoritmo.

\subsubsection{Uso de Pilha (\textit{Stack})}
\label{sec:res-stack}

A Tabela~\ref{tab:stack_kem} apresenta o uso de \textit{stack}
para os mecanismos de encapsulamento de chave (KEMs), e a
Tabela~\ref{tab:stack_sig} para os esquemas de assinatura digital.

\begin{table}[ht]
\centering
\caption{Uso de memória \textit{stack} por KEM e operação (bytes).
Medido via \textit{stack painting} com \texttt{pthread} de
4\,MB e compilação \texttt{-O0}.}
\label{tab:stack_kem}
\resizebox{\columnwidth}{!}{%
\begin{tabular}{lrrr}
\toprule
\textbf{Algoritmo} & \textbf{Keygen} & \textbf{Encap} & \textbf{Decap} \\
\midrule
\multicolumn{4}{l}{\textit{ML-KEM (FIPS\,203)}} \\[2pt]
\quad ML-KEM-512  &  13.520 &  16.160 &  16.960 \\
\quad ML-KEM-768  &  17.648 &  20.784 &  21.904 \\
\quad ML-KEM-1024 &  22.768 &  26.448 &  28.048 \\
\midrule
\multicolumn{4}{l}{\textit{HQC}} \\[2pt]
\quad HQC-128 &  38.576 &  56.448 &  60.896 \\
\quad HQC-192 &  70.688 & 106.752 & 115.760 \\
\quad HQC-256 & 108.768 & 166.592 & 181.040 \\
\bottomrule
\end{tabular}}
\end{table}

\begin{table}[ht]
\centering
\caption{Uso de memória \textit{stack} por esquema de assinatura
digital e operação (bytes). Medido via \textit{stack painting}
com \texttt{pthread} de 4\,MB e compilação \texttt{-O0}.
Medição pendente para variantes SLH-DSA-SHAKE.}
\label{tab:stack_sig}
\resizebox{\columnwidth}{!}{%
\begin{tabular}{lrrr}
\toprule
\textbf{Algoritmo} & \textbf{Keygen} & \textbf{Sign} & \textbf{Verify} \\
\midrule
\multicolumn{4}{l}{\textit{ML-DSA (FIPS\,204)}} \\[2pt]
\quad ML-DSA-44 &  28.304 &  53.456 &  25.968 \\
\quad ML-DSA-65 &  33.456 &  83.312 &  29.072 \\
\quad ML-DSA-87 &  41.648 & 117.072 &  35.472 \\
\midrule
\multicolumn{4}{l}{\textit{FN-DSA (FIPS\,206)}} \\[2pt]
\quad FN-DSA-512  &  23.608 &  49.584 &   9.064 \\
\quad FN-DSA-1024 &  40.504 &  90.736 &  13.160 \\
\midrule
\multicolumn{4}{l}{\textit{SLH-DSA-SHAKE (FIPS\,205)}} \\[2pt]
\quad SLH-DSA-SHAKE-128s & --- & --- & --- \\
\quad SLH-DSA-SHAKE-128f & --- & --- & --- \\
\quad SLH-DSA-SHAKE-192s & --- & --- & --- \\
\quad SLH-DSA-SHAKE-192f & --- & --- & --- \\
\quad SLH-DSA-SHAKE-256s & --- & --- & --- \\
\quad SLH-DSA-SHAKE-256f & --- & --- & --- \\
\bottomrule
\end{tabular}}
\end{table}

O uso de \textit{stack} é a métrica que melhor diferencia os
algoritmos sob análise.
Entre os KEMs, o ML-KEM apresenta consumo reduzido e previsível:
a operação mais custosa, \textit{decap} do ML-KEM-1024, consome
apenas 28.048 bytes.
O HQC, por outro lado, exige consideravelmente mais pilha: o
\textit{decap} do HQC-256 chega a 181.040 bytes, cerca de
6,5 vezes o valor equivalente do ML-KEM-1024.
Esse comportamento é esperado dado o tamanho dos vetores
polinomiais e os buffers intermediários que o HQC emprega
internamente.

Entre os esquemas de assinatura, o ML-DSA é o mais exigente:
a operação \textit{sign} do ML-DSA-87 consome 117.072 bytes,
reflexo dos vetores de comprometimento e das matrizes
de desafio manipulados durante a geração da assinatura.
O FN-DSA (Falcon) destaca-se pelo baixíssimo consumo de
\textit{stack} na verificação: apenas 9.064 bytes para FN-DSA-512
e 13.160 bytes para FN-DSA-1024, tornando-o especialmente atraente
para ambientes com restrições de pilha.
A geração de chaves e assinatura do FN-DSA também são competitivas
frente ao ML-DSA no mesmo nível de segurança.

As medições de \textit{stack} para as variantes SLH-DSA-SHAKE
(FIPS~205) foram adicionadas ao escopo após a coleta inicial e
não estão disponíveis para esta métrica; os dados de
\textit{heap} dessas variantes estão, contudo, presentes na
Seção~\ref{sec:res-heap}.

\subsubsection{Pico de Heap (Valgrind \texttt{massif})}
\label{sec:res-heap}

As Tabelas~\ref{tab:heap_kem} e~\ref{tab:heap_sig} apresentam o
pico de uso de \textit{heap} medido pelo
\texttt{valgrind --tool=massif}.

\begin{table}[ht]
\centering
\caption{Pico de \textit{heap} por KEM e operação (KB).
Medido com \texttt{valgrind --tool=massif}; valor de pico
de \texttt{mem\_heap\_B} para cada operação isolada.}
\label{tab:heap_kem}
\resizebox{\columnwidth}{!}{%
\begin{tabular}{lrrr}
\toprule
\textbf{Algoritmo} & \textbf{Keygen (KB)} & \textbf{Encap (KB)} & \textbf{Decap (KB)} \\
\midrule
\multicolumn{4}{l}{\textit{ML-KEM (FIPS\,203)}} \\[2pt]
\quad ML-KEM-512  & 204,97 & 204,97 & 204,97 \\
\quad ML-KEM-768  & 207,17 & 207,17 & 207,17 \\
\quad ML-KEM-1024 & 208,00 & 208,17 & 208,17 \\
\midrule
\multicolumn{4}{l}{\textit{HQC}} \\[2pt]
\quad HQC-128 & 211,07 & 211,07 & 211,07 \\
\quad HQC-192 & 219,74 & 219,74 & 219,74 \\
\quad HQC-256 & 230,60 & 230,60 & 230,60 \\
\bottomrule
\end{tabular}}
\end{table}

\begin{table}[ht]
\centering
\caption{Pico de \textit{heap} por esquema de assinatura
digital e operação (KB).
Medido com \texttt{valgrind --tool=massif}; valor de pico
de \texttt{mem\_heap\_B} para cada operação isolada.}
\label{tab:heap_sig}
\resizebox{\columnwidth}{!}{%
\begin{tabular}{lrrr}
\toprule
\textbf{Algoritmo} & \textbf{Keygen (KB)} & \textbf{Sign (KB)} & \textbf{Verify (KB)} \\
\midrule
\multicolumn{4}{l}{\textit{ML-DSA (FIPS\,204)}} \\[2pt]
\quad ML-DSA-44 & 207,92 & 208,90 & 208,90 \\
\quad ML-DSA-65 & 211,02 & 211,02 & 211,83 \\
\quad ML-DSA-87 & 214,37 & 214,59 & 214,59 \\
\midrule
\multicolumn{4}{l}{\textit{FN-DSA (FIPS\,206)}} \\[2pt]
\quad FN-DSA-512  & 204,64 & 204,81 & 204,81 \\
\quad FN-DSA-1024 & 207,16 & 207,38 & 207,38 \\
\midrule
\multicolumn{4}{l}{\textit{SLH-DSA-SHAKE (FIPS\,205)}} \\[2pt]
\quad SLH-DSA-SHAKE-128s & 210,31 & 210,31 & 210,31 \\
\quad SLH-DSA-SHAKE-128f & 219,32 & 219,32 & 219,32 \\
\quad SLH-DSA-SHAKE-192s & 218,52 & 218,52 & 218,52 \\
\quad SLH-DSA-SHAKE-192f & 237,51 & 237,51 & 237,51 \\
\quad SLH-DSA-SHAKE-256s & 231,82 & 231,82 & 231,82 \\
\quad SLH-DSA-SHAKE-256f & 251,42 & 251,42 & 251,42 \\
\bottomrule
\end{tabular}}
\end{table}

Ao contrário do \textit{stack}, o pico de \textit{heap} não
constitui um fator diferenciador expressivo entre os algoritmos.
Os valores observados situam-se em uma faixa estreita, de
aproximadamente 204,64\,KB (FN-DSA-512) a 251,42\,KB
(SLH-DSA-SHAKE-256f), uma variação total inferior a 47\,KB
em todo o conjunto avaliado.

Essa uniformidade é explicada pela arquitetura interna do
\texttt{liboqs}: a biblioteca realiza uma alocação de contexto
durante a inicialização do objeto de algoritmo
(\texttt{OQS\_KEM\_new} / \texttt{OQS\_SIG\_new}), que
representa o grosso do \textit{heap} medido pelo Valgrind,
independentemente do algoritmo selecionado.
As alocações transitórias da operação em si — vetores de chave,
texto cifrado, assinatura — são adicionadas durante a execução
e liberadas ao término, resultando em variações pequenas
entre operações do mesmo algoritmo.

O crescimento do pico de \textit{heap} entre variantes de maior
nível de segurança reflete principalmente o tamanho dos buffers
internos: o SLH-DSA-SHAKE-256f, com parâmetros mais conservadores
e assinaturas maiores, apresenta o maior pico (251,42\,KB), enquanto
o FN-DSA-512 registra o menor (204,64\,KB).
Para efeitos de seleção de algoritmo em sistemas com restrição de
memória dinâmica, contudo, o \textit{heap} não é o fator
limitante — todos os algoritmos operam dentro de uma faixa
compatível com sistemas de propósito geral.


% ============================================================
% BLOCO 8 — Seção 6 (Conclusão)
%
% LOCALIZAÇÃO: parágrafo que começa com (ou contém):
%   "A análise de consumo de memória não foi realizada,
%    pois requer uma infraestrutura de perfilamento separada
%    (como Valgrind --tool=massif)"
% AÇÃO: REMOVER esse parágrafo inteiro e substituir pelo
%       trecho abaixo, que reporta os achados reais.
% ============================================================

% --- SUBSTITUIR o parágrafo sobre memória na Conclusão ---
%
% (Remova o parágrafo original que diz que a análise não foi
%  realizada e coloque este em seu lugar, dentro do fluxo
%  das conclusões.)

O consumo de memória revelou que o \textit{stack} é a dimensão
mais discriminante entre os algoritmos.
O HQC exige até 6,5 vezes mais pilha que o ML-KEM no nível de
segurança equivalente, reflexo de seu design baseado em vetores
densos.
O FN-DSA (Falcon) destacou-se pelo consumo de \textit{stack}
excepcionalmente baixo na verificação (abaixo de 14\,KB),
sendo o candidato mais adequado para ambientes com restrições
severas de pilha.
O pico de \textit{heap}, por sua vez, é dominado pelo
\textit{overhead} de inicialização do \texttt{liboqs} e varia
menos de 47\,KB entre todos os algoritmos avaliados, não
constituindo critério diferenciador para seleção em sistemas
de propósito geral.


% ============================================================
% BLOCO 9 — Seção 6 (Trabalhos Futuros)
%
% LOCALIZAÇÃO: parágrafo de Trabalhos Futuros que começa com
%   "\textbf{Perfilamento de consumo de memória.}"
%   (ou texto equivalente propondo a análise de memória
%    como trabalho futuro)
% AÇÃO: REMOVER esse parágrafo inteiro.
%   A análise de memória foi realizada e seus resultados
%   estão na Seção~\ref{sec:res-memoria}.
%   Nenhum texto substituto é necessário neste local.
% ============================================================

% (Sem conteúdo — apenas remova o parágrafo mencionado acima.)


% ============================================================
% FIM DO ARQUIVO
% ============================================================
